% concatenated from Banerjee.tex
\documentclass[11pt,reqno]{amsart}
\usepackage{amssymb, amsmath, amsthm, mathrsfs, setspace, enumerate}
\usepackage[bookmarksnumbered, colorlinks, plainpages]{hyperref}
\usepackage{array}
\usepackage{tabularx}
\usepackage{tikz}
\usepackage{mathrsfs}
\textheight 22.78truecm 
\textwidth 15.5truecm
\setlength{\oddsidemargin}{0.35in}
\setlength{\evensidemargin}{0.35in}
\setlength{\topmargin}{-.5cm}
\newtheorem*{theo3A}{Theorem 3.A}
\newtheorem*{theo3B}{Theorem 3.B}
\newtheorem*{theo3C}{Theorem 3.C}
\newtheorem*{theo3D}{Theorem 3.D}
\newtheorem*{theo3E}{Theorem 3.E}
\newtheorem*{theo3F}{Theorem 3.F}
\newtheorem*{theo2G}{Theorem 2.G}
\newtheorem*{theo2H}{Theorem 2.H}
\newtheorem*{theo2I}{Theorem 2.I}
\newtheorem*{theo2J}{Theorem 2.J}
\newtheorem*{theo2K}{Theorem 2.K}
\newtheorem*{lem2A}{Lemma 2.A}
\newtheorem*{lem2B}{Lemma 2.B}

\newtheorem*{propo1A}{Proposition 1.A}
\newtheorem*{propo1B}{Proposition 1.B}
\newtheorem*{propo1C}{Proposition 1.C}

\newtheorem*{defi1A}{Definition 1.A}
\newtheorem*{defi1B}{Definition 1.B}
\newtheorem*{defi1C}{Definition 1.C}

\newtheorem*{defi1D}{Definition 1.D}
\newtheorem*{defi1E}{Definition 1.E}
\newtheorem*{defi1F}{Definition 1.F}
\newtheorem*{defi1G}{Definition 1.G}
\newtheorem*{defi1H}{Definition 1.H}
\newtheorem*{defi1I}{Definition 1.I}

\newtheorem*{ques2A}{Question 2.A}
\newtheorem*{ques2B}{Question 2.B}

\newtheorem{ques}{Question}[section]

\newtheorem*{exm3A}{Example 3.A}
\newtheorem*{exm3B}{Example 3.B}
\newtheorem*{exm3C}{Example 3.C}
\newtheorem*{exm3D}{Example 3.D}
\newtheorem{prop}{Proposition}[section]
\newtheorem{theo}{Theorem}[section]
\newtheorem{lem}{Lemma}[section]
\newtheorem{cor}{Corollary}[section]
\newtheorem{note}{Note}[section]
\newtheorem{exm}{Example}[section]
\newtheorem{defi}{Definition}[section]
\newtheorem{rem}{Remark}[section]
\newtheorem{propo}{Proposition}[section]
%\renewcommand{\vline}{\mid}

\newcommand{\ol}{\overline}
\newcommand{\ul}{\underline}

\newcommand{\be}{\begin{equation}}
\newcommand{\ee}{\end{equation}}
\newcommand{\beas}{\begin{eqnarray*}}
\newcommand{\eeas}{\end{eqnarray*}}
\newcommand{\bea}{\begin{eqnarray}}
\newcommand{\eea}{\end{eqnarray}}
\newcommand{\lra}{\longrightarrow}
\numberwithin{equation}{section}
\begin{document}

\title[B\MakeLowercase{orel Exceptional Values in Several Complex Variables}......]{\LARGE B\Large\MakeLowercase{orel Exceptional Values in Several Complex Variables\newline and Their Applications to Shared Values of Shifts\newline \vspace {.25cc} and Difference Operators}}

\date{}
\author[A. B\MakeLowercase{anerjee}, S. M\MakeLowercase{ajumder}, J. B\MakeLowercase{anerjee}]{A\MakeLowercase{bhijit} B\MakeLowercase{anerjee}$^1$, S\MakeLowercase{ujoy} M\MakeLowercase{ajumder}$^2$ $^{*}$, J\MakeLowercase{hilik} B\MakeLowercase{anerjee}$^3$}
\address{$^{1}$ Department of Mathematics, University of Kalyani, West Bengal 741235, India.}
\email{abanerjee\_kal@yahoo.co.in}
\address{$^{2}$Department of Mathematics, Raiganj University, Raiganj, West Bengal-733134, India.}
\email{sm05math@gmail.com, sjm@raiganjuniversity.ac.in}
	\address{$^{3}$ Department of Mathematics, University of Kalyani, West Bengal 741235, India.}
\email{jhilikbanerjee38@gmail.com}


\renewcommand{\thefootnote}{}
\footnote{2020 \emph{Mathematics Subject Classification}: 32A20, 32A22, 32H30.}
\footnote{\emph{Key words and phrases}: Meromorphic functions in $\mathbb{C}^n$, Exponent of convergence, Borel exceptional value, Order, Deficiency of values, Uniqueness, Difference Operator.}
\footnote{*\emph{Corresponding Author}: Sujoy Majumder.}

\renewcommand{\thefootnote}{\arabic{footnote}}
\setcounter{footnote}{0}

\begin{abstract} 
In this paper, we investigate shared value problems for shifts and higher-order difference operators of meromorphic and entire functions in several complex variables. Using Nevanlinna theory in $\mathbb{C}^n$, we obtain new uniqueness theorems when functions share values counting or ignoring multiplicities, extending several classical one-variable results to higher dimensions. A key contribution of this work appears in Section 2, where we establish fundamental results on Borel exceptional values in several complex variables. These propositions provide the main tools for proving our principal theorems. As applications, we derive conditions ensuring that a transcendental entire function satisfies $\Delta_c^{k} \equiv d\hspace{.05cc} f$ and we study meromorphic solutions of certain partial differential-difference equations, obtaining growth estimates and structural descriptions of entire solutions. To the best of our knowledge, this is the first systematic study of such shared value problems for higher-order difference operators in several complex variables.
\end{abstract}

\thanks{Typeset by \AmS -\LaTeX}
\maketitle


\section{{\bf Preliminaries on Nevanlinna Theory in Several Complex Variables}}

This section collects the fundamental concepts and notation from Nevanlinna theory in several complex variables that will be used throughout the paper. Our presentation follows the framework developed by Stoll \cite{Stoll-1974}, Nishino \cite{Nishino-2001}, and Hu-Li-Yang \cite{Hu-Li-Yang-2003}.

\subsection*{Basic Notation}

For any set $S$, we denote by $S^n$ its $n$-fold Cartesian product.  
If $S$ is partially ordered and $a \le b$, we use
\[
S[a,b]=\{x\in S:a\le x\le b\}, \quad
S(a,b]=\{x\in S:a<x\le b\},
\]
\[
S[a,b)=\{x\in S:a\le x<b\}, \quad
S(a,b)=\{x\in S:a<x<b\}.
\]
In particular,
\[
\mathbb{Z}_+ = \mathbb{Z}[0,\infty), \quad
\mathbb{R}_+ = \mathbb{R}[0,\infty), \quad
\mathbb{Z}^+ = \mathbb{Z}(0,\infty), \quad
\mathbb{R}^+ = \mathbb{R}(0,\infty).
\]

\subsection*{Hermitian Structure and Differential Forms}

Let $W$ be an $n$-dimensional Hermitian vector space equipped with a positive definite Hermitian form $(\cdot \mid \cdot)$. The induced norm is
\[
\|z\| = \sqrt{(z \mid z)}, \quad z \in W.
\]
When $W=\mathbb{C}^n$, the Hermitian product is taken to be
\[
(a \mid b)=\sum_{j=1}^n a_j \overline{b_j},
\quad a,b\in\mathbb{C}^n.
\]

For any subset $A\subset W$ and $r\ge0$, define (see \cite[pp. 6]{Stoll-1974})
\[
A[r]=\{z\in A:\|z\|\le r\},\quad
A(r)=\{z\in A:\|z\|<r\},\quad
A\langle r\rangle=\{z\in A:\|z\|=r\}.
\]

We introduce the functions
\[
\tau_0(z)=\|z\|,\qquad \tau_W(z)=\|z\|^2.
\]

On $W$, the exterior derivative decomposes as $d=\partial+\bar{\partial}$, and we set
\[
d^c=\frac{\iota}{4\pi}(\bar{\partial}-\partial), \qquad
dd^c=\frac{\iota}{2\pi}\partial\bar{\partial}.
\]
The standard \text{K\"ahler} form on $W$ is
\[
\upsilon_W = dd^c \tau_W >0.
\]
On $W\setminus\{0\}$, we further define
\[
\omega_W = dd^c \log \tau_W \ge 0, \qquad
\sigma_W = d^c\log\tau_W \wedge \omega_W^{n-1},
\]
where $n=\dim(W)$ (see \cite[pp. 6]{Stoll-1974}).

\subsection*{Projective Space}

\begin{defi1A}\cite[pp. 7]{Nishino-2001}
	Let $z',z''\in\mathbb{C}^{n+1}\setminus\{0\}$. We say that $z'$ and $z''$ are equivalent if there exists $c\in\mathbb{C}\setminus\{0\}$ such that
	\[
	z_j''=c z_j', \quad j=0,1,\ldots,n.
	\]
	The equivalence classes in $\mathbb{C}^{n+1} \backslash  \{0\}$ form an 
$n$-dimensional space called the complex projective space, which we denote by
	\[
	\mathbb{P}^n=\mathbb{P}(\mathbb{C}^{n+1}).
	\]
	Identify the compactified plane $\mathbb{C}\cup\{\infty\}$ with $\mathbb{P}^1=\mathbb{P}(\mathbb{C}^2)$ by setting (see \cite[pp. 634-635]{Stoll-1983})
	\[
	\mathbb{P}(z,w)=\frac{z}{w}, \quad
	\mathbb{P}(z,1)=z, \quad
	\mathbb{P}(1,0)=\infty
	\]
	if $z\in\mathbb{C}$ and $w\in\mathbb{C}\backslash \{0\}$.
\end{defi1A}

\subsection*{Divisors and Multiplicities}

\begin{defi1B}\cite[pp. 12]{Stoll-1974}
	Let $G\neq \varnothing$ be an open subset of $W$. Let $f$ be a holomorphic function in $W$.
Take $a\in G$. Let $G_a$ be the connectivity component of $G$ containing $a$. Assume $f\mid_{G_a}\not\equiv 0$. Then $f$ admits a local expansion
	\[
	f(z)=\sum_{\lambda=p}^{\infty}P_\lambda(z-a),
	\]
	where $P_\lambda$ is homogeneous of degree $\lambda$ and $P_p\not\equiv0$. The polynomials $P_{\lambda}$ depend on $f$ and $a$ only.
	The integer $\mu_f^0(a)=p$ is called the zero multiplicity of $f$ at $a$.
\end{defi1B}

\begin{defi1C}\cite[pp. 12]{Stoll-1974}
Let $f$ be a meromorphic function on $G$, where $G\neq \varnothing$ is an open subset of $W$.
Take $a\in G$ and $c\in\mathbb{P}^1$. Let $G_a$ be the component of $G$ containing $a$. If $0\equiv f\mid_{G_a}\not\equiv c$, define $\mu^c_f(a)=0$. Assume $0\not\equiv f\mid_{G_a}\not\equiv c$. Then an open connected neighborhood $U$ of $a$ in $G$ and holomorphic functions $g\not\equiv 0$ and $h\not\equiv 0$ exist on $U$ such that $h. f\mid_U=g$ and $\dim g^{-1}(0)\cap h^{-1}(0)\leq n-2$, where $n=\dim(W)$. Therefore the $c$-multiplicity of $f$ is just $\mu^c_f=\mu^0_{g-ch}$ if $c\in\mathbb{C}$ and $\mu^c_f=\mu^0_h$ if $c=\infty$. The function $\mu^c_f:G\to \mathbb{Z}$ is nonnegative and is called the $c$-divisor of $f$. 
\end{defi1C}

If $f\not\equiv0$ on each component of $G$, the divisor of $f$ is
\[
\mu_f=\mu_f^0-\mu_f^\infty.
\]
The function $f$ is holomorphic if and only if $\mu_f\ge0$.

\begin{defi1D}\cite[pp. 13]{Stoll-1974}
A function $\nu:G\to \mathbb{Z}$ is said to be a divisor if and only if for each $a\in G$ an open, connected neighborhood $U$ of $a$ in $G$ and a meromorphic function $f\not\equiv 0$ exist such that $\nu\mid_{U}=\mu_f$. A divisor $\nu:G\to \mathbb{Z}$ is non-negative if and only if for each $a\in G$ an open, connected neighborhood $U$ of $a$ in $G$ and a holomorphic function $f\not\equiv 0$ exist such that $\nu\mid_{U}=\mu_f$.
\end{defi1D}

We denote
\[
\operatorname{supp}\nu=\overline{\{z\in G:\nu(z)\ne0\}}.
\]


\begin{defi1E}\cite[pp. 26]{Stoll-1974} Let $\nu\geq 0$ be a non-negative divisor on the hermitian vector space $W$. Define $A=\displaystyle \operatorname{supp}\nu$. Assume $0\not\in A$. A non-negative, increasing, integral valued function $q:\mathbb{R}^+\to \mathbb{Z}^+$ is said to be a weight for $\nu$ if and only if the integral
	\begin{align*}
		\int\limits_{A} \nu\;\left(r/\tau_0\right)^{q\circ\tau_0+1}\omega_W^{n-1}<\infty
	\end{align*}
	converges for all $r>0$, where $n=\dim(W)$. For every given non-negative divisor $\nu$ there exists a weight function for $\nu$ (see \cite{Stoll-1953}).
\end{defi1E}
\subsection*{Counting and Characteristic Functions}

\begin{defi1F} Take $0<R\leq +\infty$. Let $\nu$ be a divisor on $W(R)$ with $A=\displaystyle \operatorname{supp}\nu$. For $t>0$, the counting function $n_{\nu}$ is defined by
\begin{align*}
 \displaystyle n_{\nu}(t)=t^{-2(m-1)}\int_{A[t]}\nu\; \upsilon_W^{n-1},
 \end{align*}
where $n=\dim(W)$.
Obviously 
\begin{align*}
n_{\nu}(0)=\lim\limits_{t\to 0}n_{\nu}(t)\in\mathbb{Z}
\end{align*}
exists, where the Lelong number $n_{\nu}(0)$ is an integer (see \cite[Theorem 3.1]{Stoll-1974}). We know that
\begin{align*}
 \displaystyle n_{\nu}(t)=\int_{A(t)}\nu\;\omega_W^{n-1}+n_{\nu}(0).
 \end{align*}
If $\nu$ is non-negative, then $n_{\nu}$ increases. For $0<s<r<R$, we define the valence function of $\nu$ by
\begin{align*}
N_{\nu}(r)=N_{\nu}(r,s)=\int_{s}^r n_{\nu}(t)\frac{dt}{t}.
\end{align*}
\end{defi1F}

For $\nu=\mu_f^a$, we write $n(t,a;f)$ and $N(r,a;f)$, with the usual conventions for $a=\infty$.  
Truncated multiplicities are defined by $\mu_{f,k}^a=\min\{\mu_f^a,k\}$.
\vspace{1.2mm}

For $k\in\mathbb{N}$, define the truncated multiplicity functions on $W$ by $\mu_{f,k}^a(z)=\min\{\mu_f^a(z),k\}$.
We write the truncated counting functions 
\beas n_{\nu}(t)=\begin{cases}
n_k(t,a;f), &\text{if $\nu=\mu_{f,k}^a$}\\
\ol n(t,a;f), &\text{if $\nu=\mu_{f,1}^a$}
\end{cases}\eeas
and the truncated valence functions
\beas N_{\nu}(t)=\begin{cases}
N_k(t,a;f), &\text{if $\nu=\mu_{f,k}^a$}\\
\ol N(t,a;f), &\text{if $\nu=\mu_{f,1}^a$}.
\end{cases}\eeas

If $a=\infty$, we write $n_k(t,a;f)=n_k(t,f)$ and $N_k(t,a;f)=N_k(t,f)$
and so on.\par
A non-negative divisor $\nu:W\to \mathbb{Z}_+$ is said to be algebraic if and only if $\nu$ is  the zero divisor of a polynomial. Thus a divisor $\nu:W\to \mathbb{Z}_+$ is algebraic if and only if $n_{\nu}$ is bounded, which implies that $N_{\nu}=O(\log r)$ (see \cite[pp. 19]{Stoll-1974}).

\subsection*{Nevanlinna Characteristic}

\begin{defi1G}\cite[pp. 16-17]{Stoll-1974}
Let $G\neq \varnothing$ be an open subset in $W$. Let $f$ be a meromorphic function in $G$ in the sense that $f$ can be written as a quotient of two relatively prime holomorphic functions. We will write $f = (f_0, f_1)$ where $f_0 \not\equiv 0$. The standard definition of the Nevanlinna characteristic function of $f$ is given by
\begin{align*}
T_f(r,s) := \int_s^r \frac{A_f(t)}{t}\, dt,
\end{align*}
where $0 < s < r$ and
\begin{align*}
A_f(t)
= \frac{1}{t^{2n-2}} \int_{W(t)} f^*(\ddot{\omega}) \wedge \upsilon_W^{n-1}
= \int_{W(t)} f^*(\ddot{\omega}) \wedge \omega_W^{\,n-1} + A_f(0),
\end{align*}
where $n=\dim(W)$.
Here the pullback $f^*(\ddot{\omega})$ satisfies
\begin{align*}
f^*(\ddot{\omega}) = dd^c \log\left( |f_0|^2 + |f_1|^2 \right)
\end{align*}
for all $z$ outside of the set of indeterminacy $I_f:= \{ z \in W : f_0(z) = f_1(z) = 0 \}$ of $f$.

Take $a\in \mathbb{P}^1$ and $0<R\leq +\infty$. Let $f\not\equiv 0$ be a meromorphic function on $W(R)$.
For $0<r<R$, define the compensation of $f$ for $a$ by
\begin{align*}
m^a_f(r)=\int\limits_{W\langle r\rangle}\log \frac{1}{||f,a||} \;\sigma_W,
\end{align*}
where $||f,a||$ denotes the chordal distance from $f$ to $a\in\mathbb{P}^1$.
Then the First Main Theorem of Nevanlinna theory becomes
\begin{align*}
T_f(r)=T_f(r,s)=N_{\mu^a_f}(r,s)+m^a_f(r)-m^a_f(s),
\end{align*}
where $0<s<r$.	
\end{defi1G}

There is slightly different way to continue the formulation of Nevanlinna theory from here (see \cite[pp.15]{Hu-Li-Yang-2003}).
Take $0<R\leq +\infty$. Let $f\not\equiv 0$ be a meromorphic function on $W(R)$. Let $0<s<r<R$. 
Now with the help of the positive logarithm function, we define the proximity function of $f$ by
\begin{align*}
\displaystyle m(r, f)=\int_{W\langle r\rangle} \log^+ |f|\;\sigma_W \geq  0.
\end{align*}

The characteristic function of $f$ is defined by $T(r,f)=m(r,f)+N(r,f)$. We know that (see \cite[pp.15]{Hu-Li-Yang-2003})
\begin{align*}
\displaystyle T\left(r,\frac{1}{f}\right)=T(r,f)-\int_{W\langle s\rangle} \log |f|\;\sigma_W.
\end{align*}

We define $m(r,a;f)=m(r,f)$ if $a=\infty$ and $m(r,a;f)=m\big(r,\frac{1}{f-a}\big)$ if $a$ is finite complex number. Now if $a\in\mathbb{C}$, then the first main theorem of Nevanlinna theory becomes $m(r,a;f)+N(r,a;f)=T(r,f) + O(1)$, where $O(1)$ denotes a bounded function when $r$ is sufficiently large.

Finally, if we compare the functions $T_f(r)$ and $T(r,f)$, then we have (see \cite[pp.19]{Hu-Li-Yang-2003})
\begin{align*}
T_f(r)=T(r,f)+O(1).
\end{align*}


\subsection*{Growth, Exponent of convergence and Deficiency}
Consider an increasing, non-negative function $S:\mathbb{R}_{+}\to\mathbb{R}_{+}.$
The order of $S$ is defined by (see \cite[pp. 28]{Stoll-1974})
\begin{align*}
\operatorname{Ord}\, S= \limsup_{r\to\infty}\frac{\log^{+} S(r)}{\log r}.
\end{align*}

For $0<\mu\in\mathbb{R}$, we define
\begin{align*}
t_{\mu}(S)=\limsup_{r\to\infty}\frac{S(r)}{r^{\mu}}\quad \text{and}\quad
J_{\mu}(S)=\int_{1}^{\infty}\frac{S(t)}{t^{\mu+1}}\,dt.
\end{align*}

Clearly if $0<\mu<\operatorname{Ord}S$, then $t_{\mu}(S)=J_{\mu}(S)=\infty$ and if $\mu>\operatorname{Ord}S$, then $t_{\mu}(S)<\infty$ and $\ J_{\mu}(S)<\infty$.
If $0<\operatorname{Ord}S=\lambda<\infty$, the function $S$ is said to be of
\begin{itemize}
	\item maximal type \quad iff \quad \(t_{\lambda}(S)=+\infty\),
	\item middle type \quad iff \quad \(0<t_{\lambda}(S)<+\infty\),
	\item minimal type \quad iff \quad \(t_{\lambda}(S)=0\).
\end{itemize}

The function $S$ is said to be in the
\begin{itemize}
	\item convergence class \quad iff \quad \(J_{\lambda}(S)<+\infty\),
	\item divergence class \quad iff \quad \(J_{\lambda}(S)=+\infty\).
\end{itemize}


\begin{defi1H}\cite[pp. 28]{Stoll-1974}Let $\nu\geq 0$ be a non-negative divisor on $W$. Then the order, class, and type
of $\nu$ are defined as the order, class, and type of $n_{\nu}$. Observe that for each $s>0$, the order, class and type of $n_{\nu}$
coincide with the order, class and type of $N_{\nu}(\square,s)$. If $0\notin \operatorname{supp}\nu$, this is true for $s=0$ also.
\end{defi1H}

Let $f$ be a non-constant meromorphic function in $W$. The function $f$ is said to have order $\rho(f)$, maximal, mean, minimal type or convergence class if the characteristic function $T(r,f)$ has those properties. Consequently the order and the lower order of $f$ are given respectively by
\[\operatorname{Odr} f= \rho(f):=\varlimsup\limits_{r\to \infty}\frac{\log^+ T(r,f)}{\log r}\;\;\text{and}\;\;\mu(f):=\varliminf\limits_{r\to \infty}\frac{\log^+ T(r,f)}{\log r}.\]

\begin{propo1A}\cite[Proposition 1.7]{Hu-Li-Yang-2003}
Suppose that $f$ and $g$ are non-constant meromorphic functions in $\mathbb{C}^n$. We have
\begin{enumerate}
    \item[(a)] $\rho(\mathcal{L}(f)) = \rho(f)$, where $\mathcal{L}(f)$ is any M\"{o}bius transformation of $f$.
    \item[(b)] $\rho (f \cdot g) \leq \max\{\rho(f), \rho(g)\}$.
    \item[(c)] $\rho(f + g) \leq \max\{\rho(f), \rho(g)\}$.
    \item[(d)] if $\rho(f) < \lambda(g)$, then $\rho(fg) = \rho(f + g) = \rho(g)$.
    \item[(e)] $\max\{\mu(fg), \mu(f + g)\} \leq \max\{\rho(f), \mu(g)\}$.
    \item[(f)] if $\rho(f) < \mu(g)$, then $\max\{\mu(fg), \mu(f + g)\} \leq \mu(g)$.
    \item[(g)] if $\rho(f) < \mu(g)$, then $T(r, f) = o(T(r, g))$ as $r \to \infty$.
\end{enumerate}
\end{propo1A}

\begin{defi1I}\cite[pp. 28-29]{Stoll-1974}
Define $A=\operatorname{supp}\nu$ and $B=A-A(1).$ A number $\mu>0$ is an exponent of convergence of $\nu$ if one-hence all-of
the following integrals exist:
\begin{align*}
&\int_{B} \nu\,\tau_{0}^{-\mu}\;\omega_W^{n-1} < \infty\\
&\int_{B} \nu\;\tau_{0}^{\,1-m-\mu}\,v_W^{n-1} < \infty,\\
&\int_{1}^{\infty} n_{\nu}(t)\,t^{-\mu-1}\,dt < \infty,\\
&\int_{1}^{\infty} N_{\nu}(t,s)\,t^{-\mu-1}\,dt < \infty,
\quad \text{if } s>0.
\end{align*}

Otherwise, $\mu$ is said to be an exponent of divergence of $\nu$. If $\mu>\operatorname{Ord}\nu$, then $\mu$ is an exponent of convergence of $\nu$; if $0<\mu<\operatorname{Ord}\nu$, then $\mu$ is an exponent of divergence of $\nu$.
If $0<\operatorname{Ord}\nu=\mu<\infty$, then $\mu$ is an exponent of convergence of $\nu$ if and only if $\nu$ is in the convergence class.
\end{defi1I}

Now we state the following result related to canonical function of finite order.
\begin{propo1B} \cite[Theorem 6.3]{Stoll-1974} Let $\nu$ be a non-negative divisor on $W$ with $0\notin \operatorname{supp}\nu$.
Assume that $q+1\in\mathbb{N}$ is an exponent of convergence of $\nu$. Let $h$ be the canonical function of $\nu$ for the constant weight $q$ of $\nu$. Take $\theta\in\mathbb{R}$ and $r\in\mathbb{R}$ with $0<\theta<1$ and $r>0$. Then
\begin{align*}
&\log M_h(\theta r)\le 8m\,c(q)\,(1-\theta)^{-3m}\,K_q(r;n_\nu),\\
&\operatorname{Ord}\nu\le \operatorname{Ord} h\le \max\{q,\operatorname{Ord}\nu\}\le q+1.
\end{align*}
	
If $q\in\mathbb{Z}_{+}$ is the smallest non-negative integer such that $q+1$ is an exponent of convergence of $\nu$, then
\begin{align*}
q \le \operatorname{Ord}\nu = \operatorname{Ord} h \le q+1,
\end{align*}
\begin{align*}
&\operatorname{type} h = \operatorname{type}\nu\quad \text{provided } q < \operatorname{Ord} h \le q+1,\\
&\operatorname{class} h = \operatorname{class}\nu\quad \text{provided } q <\operatorname{Ord} h < q+1.
\end{align*}
\end{propo1B}

Proposition 1.B was proved in \cite{Stoll-1953}. Let $\nu\geq 0$ be a non-negative divisor of finite order on $W$. Then the smallest non-negative integer $q$ such that $q+1$ is an exponent of convergence of $\nu$ is called the \emph{genus} of $\nu$.

The canonical function of a non-negative divisor $\nu$ on $W$ for a constant weight $q$ can be intrinsically characterized (see \cite[Proposition 6.4]{Stoll-1974}).
\begin{propo1C}\cite[Theorem 3]{Ronkin-1968} Let $\nu\geq 0$ be a non-negative divisor on $W$ with $0\notin \operatorname{supp}\nu$.
Let $q+1\in\mathbb{N}$ be an exponent of convergence of $\nu$. Then there exists one and only one holomorphic function
$h:W\to\mathbb{C}$ such that
\begin{align*}
h(0)=1,\quad\mu_{\log h}(0)\ge q+1,\quad \mu_h=\nu,\quad\operatorname{Ord} h \le q+1
\end{align*}
and such that if $\operatorname{Ord} h=q+1$, then $h$ is of minimal type. The unique function is the canonical function of $\nu$ for the constant weight $q$.
\end{propo1C}

Clearly $f$ is non-constant, then $T(r, f) \rightarrow \infty$ as $r \rightarrow$ $\infty$. Further $f$ is rational if and only if $T(r,f)=O(\log r)$ (see \cite[pp. 19]{Stoll-1974}).
On the other hand, if $f$ is transcendental, then 
\[\lim\limits_{r \rightarrow \infty} \frac{T(r, f)}{\log r}=+\infty.\] 
The deficiency of $f$ for a value $a$ is defined us $\delta_f(a)$:
\begin{align*}
0\leq \delta_f(a)=1-\limsup\limits_{r\to \infty}\frac{N(r,a;f)}{T(r,f)}\leq 1.
\end{align*}
Clearly $\delta_f(a)=1$ if $\mu^a_f$ is algebraic and if $f$ is transcendental. The set $\{a\in\mathbb{P}^1:\delta_f(a)>0\}$ is at most countable and (see \cite[pp. 20]{Stoll-1974})
\begin{align*}
\sum\limits_{a\in\mathbb{P}^1} \delta_f(a)\leq 2.
\end{align*}
An element $a\in \mathbb{P}^1$ is said to be a Picard (exceptional) value of $f$ if $a\not\in f(W)$. Obviously, $\delta_f(a)=1$ if $a$ is a Picard value of $f$.


\section{{\bf Borel exceptional value and related results in several complex variables}}

This section constitutes the mathematical core of the paper. It establishes the main structural results on Borel exceptional values in several complex variables, presenting them in the form of clear and precise propositions. These propositions provide the essential groundwork and form the basis for the principal theorems proved in the subsequent sections. 
 
\begin{defi} Let $f$ be a transcendental meromorphic function in $\mathbb{C}^n$ with order $\rho(f)>0$. Let $a\in\mathbb{P}^1$. Then $a$ is said to be a Borel exceptional value of $f$ if $\operatorname{Ord}\mu^a_f<\rho(f)$. According to Definition 1.H, we have
\begin{align*}
\limsup\limits_{r\to \infty} \frac{\log ^+ N(r,a;f)}{\log r}<\rho(f).
\end{align*}
\end{defi}

The following result is known as second main theorem:
\begin{lem2A}\label{L.2}\cite[Lemma 1.2]{Hu-Yang-1996} Let $f$ be a non-constant meromorphic function in $\mathbb{C}^n$ and let $a_1,a_2,\ldots,a_q$ be different points in $\mathbb{P}^1$. Then
\begin{align*}
(q-2)T(r,f)\leq \sideset{}{_{j=1}^{q}}{\sum} \ol N(r,a_j;f)+O(\log (rT(r,f)))
\end{align*}
holds only outside a set of finite measure on $\mathbb{R}^+$.
\end{lem2A}

\begin{propo}\label{pro0} Let $f$ be a transcendental meromorphic function in $\mathbb{C}^n$ with finite positive order $\rho(f)$. Then $f$ has at most two Borel exceptional values.
\end{propo}
\begin{proof} If possible suppose there are at least three distinct Borel exceptional values $a_1$, $a_2$ and $a_3$, say. Let 
\begin{align}\label{pro1:1.1}
 \limsup\limits_{r\to \infty} \frac{\log ^+ N(r,a_j;f)}{\log r}=s_j
 \end{align}
and $s=\max\{s_1,s_2,s_3\}$. Clearly $s<\rho(f)$. Now using Lemma 2.A, we get
\begin{align}\label{pro1:1.2}
T(r,f)\leq \sideset{}{_{j=1}^{3}}{\sum} \ol N(r,a_j;f)+O(\log (rT(r,f)))
\end{align}
holds only outside a set of finite measure on $\mathbb{R}^+$. Taking $\log^+$ on the both sides of (\ref{pro1:1.2}) and then using (\ref{pro1:1.1}), we get
\begin{align*}
\rho(f)=\limsup\limits_{r\to \infty} \frac{\log ^+ T(r,f)}{\log r}\leq s<\rho(f),
\end{align*}
which is impossible. Hence $f$ has at most two Borel exceptional values.
\end{proof}

\begin{rem} Obviously, when $\rho(f)>0$, Picard exceptional value is the Borel exceptional value. Then $1$ is the Borel exceptional value of $f(z)=e^{z_1+z_2+\ldots+z_n}+1$.
\end{rem}

\begin{lem2B}\label{2A}\cite[Lemma 3.59]{Hu-Li-Yang-2003} Let $P$ be a non-constant entire function in $\mathbb{C}^n$. Then 
\begin{align*}
\rho(e^P)=
\begin{cases}
\deg(P), & \text{if $P$ is a polynomial,}\\
+\infty, & \text{otherwise.}
\end{cases}
\end{align*}
\end{lem2B}

\begin{propo}\label{Pro1} Let $f$ be a non-constant meromorphic function in $\mathbb{C}^n$ such that $f(0)\neq 0,\infty$ and $\rho(f)$ be finite. Then $$
f=\frac{H_1}{H_2}e^{\alpha},$$
where $H_1$ and $H_2$ are canonical functions such that $\mu^0_f=\mu^0_{H_1}$, $\mu^{\infty}_f=\mu^0_{H_2}$,
\begin{align*}
\rho(H_1)=\limsup\limits_{r\to\infty} \frac{\log^+ N_{\mu^0_{f}}(r,0)}{\log r}\leq \rho(f)
\end{align*}
and 
\begin{align*}
\rho(H_2)=\limsup\limits_{r\to\infty} \frac{\log^+ N_{\mu^{\infty}_{f}}(r,0)}{\log r}\leq \rho(f)
\end{align*}
and $\alpha$ is a polynomial in $\mathbb{C}^n$ such that $\deg(\alpha)=\rho(e^{\alpha})\leq \rho(f)$.
\end{propo}
\begin{proof} We note that for every meromorphic function $f$ in $\mathbb{C}^n$ there are entire functions $g$ and 
$h$ such that 
\begin{align}\label{p1.1}
f=\frac{g}{h}
\end{align}
in $\mathbb{C}^n$ and $g$ and $h$ are coprime at every point of $\mathbb{C}^n$ \cite[pp. 186]{LH1} which implies 
that $F=(h,g):\mathbb{C}^n\to \mathbb{C}^2$ is a reduced representative of $f$. Since $f(0)\neq 0, \infty$, we have $g(0)\neq 0$ and $h(0)\neq 0$. Let $\nu=\mu^0_g$. Clearly $\nu$ is a non-negative divisor on $\mathbb{C}^n$ such that $0\not\in \text{supp}\;\nu$. Suppose $q\in\mathbb{Z}_+$ is the smallest non-negative integer such that $q+1$ is an exponent of convergence of $\nu$. According to Propositions 1.B and 1.C, there exists one and only one holomorphic function $H_1:\mathbb{C}^n\to \mathbb{C}$ such that $\nu=\mu^0_{H_1}$ and 
\begin{align*}
\rho(H_1)=\text{Ord}\;\nu=\limsup\limits_{r\to\infty} \frac{\log^+ N_{\mu^0_{H_1}}(r,0)}{\log r}=\limsup\limits_{r\to\infty} \frac{\log^+ N_{\mu^0_{g}}(r,0)}{\log r}=\limsup\limits_{r\to\infty} \frac{\log^+ N_{\mu^0_{f}}(r,0)}{\log r}.
\end{align*}
Now using first main theorem, we deduce that
\begin{align}\label{p1.2}
\rho(H_1)=\limsup\limits_{r\to\infty} \frac{\log^+ N_{\mu^0_{f}}(r,0)}{\log r}\leq \limsup\limits_{r\to\infty} \frac{\log^+ T(r,f)}{\log r}=\rho(f).
\end{align}
Note that $g$ and $H_1$ are entire functions in $\mathbb{C}^n$ with the same divisor $\mu^0_g=\mu^0_{H_1}$. Then there exists an entire function $\alpha_1$ in $\mathbb{C}^n$ such that $g=H_1e^{\alpha_1}$.

Again let $\nu=\mu^0_h$. Clearly $\nu$ is a non-negative divisor on $\mathbb{C}^m$ such that $0\not\in \text{supp}\;\nu$. Then there exists one and only one holomorphic function $H_2:\mathbb{C}^n\to \mathbb{C}$ such that $\nu=\mu^0_{H_2}$ and 
\begin{align}\label{p1.3}
\rho(H_2)=\limsup\limits_{r\to\infty} \frac{\log^+ N_{\mu^{\infty}_{f}}(r,0)}{\log r}\leq \limsup\limits_{r\to\infty} \frac{\log^+ T(r,f)}{\log r}=\rho(f).
\end{align}
Since $h$ and $H_2$ are entire functions in $\mathbb{C}^n$ with the same divisor $\mu^0_h=\mu^0_{H_2}$, there exists an entire function $\alpha_2$ in $\mathbb{C}^n$ such that $h=H_2e^{\alpha_2}$. If we take $\alpha=\alpha_1-\alpha_2$, then from (\ref{p1.1}), we have
\begin{align*}
f=\frac{H_1}{H_2}e^{\alpha},
\end{align*}
where $H_1$ and $H_2$ are canonical functions such that $\mu^0_f=\mu^0_{H_1}$, $\mu^{\infty}_f=\mu^0_{H_2}$ and (\ref{p1.2}) and (\ref{p1.3}) hold and $\alpha$ is an entire function in $\mathbb{C}^n$ such that $\rho(e^{\alpha})\leq \rho(f)$. Since $\rho(f)$ is finite, according to Lemma 2.B, $\alpha$ must be a polynomial in $\mathbb{C}^n$ such that $\deg(\alpha)=\rho(e^{\alpha})\leq \rho(f)$.
\end{proof}


\begin{propo}\label{Pro2} Let $h$ be a non-constant polynomial in $\mathbb{C}^n$ and let $f=e^{h}$. Then $\rho(f)=\mu(f)=\deg(h)$.
\end{propo}
\begin{proof} Note that $|h(z)|\leq C||z||^{\deg(h)}$ holds for $||z||>1$, where $C$ is a positive constant. Now for sufficiently large value of $r$, we have $T(r,f)=m\left(r,e^h\right)=O\left(r^{\deg(h)}\right)$. Therefore it is easy to verify that $\rho(f)\leq \deg(h)$ and $\mu(f)\leq \deg(h)$. Take $\xi\in\mathbb{C}^n\backslash \{0\}$. Define a holomorphic mapping $j_{\xi}:\mathbb{C}\to \mathbb{C}^n$ by $j_{\xi}(z)=z\xi$. Clearly
\begin{align*}
f_{\xi}=f\circ j_{\xi}=e^{h\circ j_{\xi}}=e^{h_{\xi}}:\mathbb{C}\to \mathbb{C}
\end{align*}
is a holomorphic mapping. Clearly $\deg(h_{\xi})=\deg(h)$. But in $\mathbb{C}$, we know that (see Theorem 1.44 \cite{Yang-Yi-2003}) $\rho(f_{\xi})=\mu(f_{\xi})=\deg(h_{\xi})=\deg(h)$. On the other hand we know that (see \cite[pp.286]{Hu-Li-Yang-2003})
\begin{align}\label{pro2:1.1} 
T_{f_{\xi}}(r,0)\leq \frac{1+\theta}{(1-\theta)^{2n-1}} T_{f}\left(\frac{r}{\theta},0\right),\quad \xi\in\mathbb{C}^n[0,1]\backslash \{0\}
\end{align}
where $0<\theta<1$. Now inequality (\ref{pro2:1.1}) implies that
$\rho(f_{\xi})\leq \rho(f)$ and $\mu(f_{\xi})\leq \mu(f)$. Consequently, we have $\deg(h)\leq \rho(f)\leq \deg(h)$ and $\deg(h)\leq \mu(f)\leq \deg(h)$. Finally, we have $\mu(f)=\rho(f)=\deg(h)$.
\end{proof}

\begin{propo}\label{Pro3} Let $f$ be a transcendental meromorphic function in $\mathbb{C}^n$ and $\rho(f)>0$ be finite. If $f$ has two distinct Borel exceptional value, say $a_1$ and $a_2$ such that $f(0)\neq a_1,a_2$ then $\rho(f)=\mu(f)\in\mathbb{N}$.
\end{propo}
\begin{proof} Without loss of generality, we assume that $a_1=0$ and $a_2=\infty$ otherwise it enough to study the function $\frac{f-a_1}{f-a_2}$. According to Proposition \ref{Pro1}, we have
\begin{align}\label{Pro3:1.1}
f=\frac{H_1}{H_2}e^{P},
\end{align}
where $H_1$ and $H_2$ are canonical functions such that
\begin{align}\label{Pro3:1.2}
\rho(H_1)=\limsup\limits_{r\to \infty}\frac{\log^+ N_{\mu^0_{H_1}}(r,0)}{\log r}=\limsup\limits_{r\to \infty}\frac{\log^+ N_{\mu^0_{f}}(r,0)}{\log r}<\rho(f),
\end{align}
\begin{align}\label{Pro3:1.3}
\rho(H_2)=\limsup\limits_{r\to \infty}\frac{\log^+ N_{\mu^0_{H_2}}(r,0)}{\log r}=\limsup\limits_{r\to \infty}\frac{\log^+ N_{\mu^{\infty}_{f}}(r,0)}{\log r}<\rho(f)
\end{align}
and $P$ is a polynomial in $\mathbb{C}^n$ such that $\deg(P)=\rho(e^P)\leq \rho(f)$. If we take 
\begin{align}\label{Pro3:1.4}
H=\frac{H_1}{H_2},
\end{align}
then from (\ref{Pro3:1.1}), we have 
\begin{align}\label{Pro3:1.5}
e^P=\frac{f}{H}.
\end{align}
Now by Proposition \ref{Pro2}, we have $\mu(e^P)=\rho(e^P)=\deg(P)$. Note that by Proposition 1.A(a), we have $\rho(1/H)=\rho(H)$.
Again by Proposition 1.A(b), we have from (\ref{Pro3:1.2})-(\ref{Pro3:1.4}) that
\begin{align*}
\rho(H)\leq \max\{\rho(H_1),\rho(H_2)\}<\rho(f).
\end{align*}
and so from (\ref{Pro3:1.5}), we have
\begin{align*}
\rho(e^P)\leq \max\{\rho(H),\rho(f)\}=\rho(f).
\end{align*}
Again from (\ref{Pro3:1.5}), we get $f=He^P$ and so by Proposition 1.A(b), we have
\begin{align*}
\rho(f)\leq \max\{\rho(e^P),\rho(H)\}=\rho(e^P).
\end{align*}
Therefore $\rho(f)=\rho(e^P)$ and so $\rho(H)<\mu(e^P)$. Since $f=He^P$, using Proposition 1.A(e), we get $\mu(f)\leq \max\{\rho(H),\mu(e^P)\}=\mu(e^P)$. 
Since $e^P=f.\frac{1}{H}$, using Proposition 1.A(e), we have 
\begin{align*}
\mu(e^P)\leq \max\{\rho(H),\mu(f)\}\leq \max\{\rho(H),\mu(e^P)\}=\mu(e^P),
\end{align*}
i.e., $\max\{\rho(H),\mu(f)\}=\mu(e^P)$. As $\rho(H)<\mu(e^P)$, it follows that $\mu(f)=\mu(e^P)=\rho(e^P)=\rho(f)$. Therefore $\mu(f)=\rho(f)=\rho(e^P)=\deg(P)\in\mathbb{N}$.
\end{proof}

\begin{propo}\label{Pro4} Let $f$ be a transcendental meromorphic function in $\mathbb{C}^n$ and $\rho(f)>0$ be finite. If $f$ has two distinct Borel exceptional value, say $a_1$ and $a_2$ such that $f(0)\neq a_1,a_2$, then $$\delta(a_1;f)=\delta(a_2;f)=1.$$
\end{propo}
\begin{proof} Without loss of generality, we assume that $a_1=0$ and $a_2=\infty$ otherwise it enough to study the function $\frac{f-a_1}{f-a_2}$. According to Proposition \ref{Pro1}, we have
\begin{align}\label{Pro4:1.1}
f=\frac{H_1}{H_2}e^{P},
\end{align}
where $H_1$ and $H_2$ are canonical functions such that (\ref{Pro3:1.2}) and (\ref{Pro3:1.3}) hold and $P$ is a polynomial in $\mathbb{C}^n$ such that $\deg(P)=\rho(e^P)\leq \rho(f)$. Now by proposition \ref{Pro3}, we have $\mu(f)=\rho(f)$. If we take $H=H_1/H_2$, then using Proposition 1.A(b), we have from (\ref{Pro3:1.2})-(\ref{Pro3:1.4}) that
\begin{align*}
\rho(H)\leq \max\{\rho(H_1),\rho(H_2)\}<\rho(f)=\mu(f).
\end{align*}
Again using Proposition 1.A(g), we conclude that $T(r,H)=o(T(r,f))$. Now in view of first main theorem and using (\ref{Pro3:1.2}) and (\ref{Pro3:1.3}), we have respectively $N(r,0;f)=o(T(r,f))$ and $N(r,\infty;f)=o(T(r,f))$. Consequently $\delta(a_1;f)=\delta(a_2;f)=1$.
\end{proof}

\begin{table}[htbp]
	\centering
	\renewcommand{\arraystretch}{1.2}
	
	\begin{tabularx}{\textwidth}{||>{\centering\arraybackslash}p{3.2cm}||
			>{\centering\arraybackslash}p{4.8cm}||
			>{\centering\arraybackslash}p{6.0cm}||}
		\hline\hline
		\textbf{Result} & \textbf{Main Tool Used} & \textbf{Purpose} \\
		\hline\hline
		Lemma~\protect\ref{L.2} & Second Main Theorem & Growth comparison \\
		\hline
		Proposition~\protect\ref{pro0} & Lemma~\protect\ref{L.2} & Maximum two exceptional values \\
		\hline
		Proposition~\protect\ref{Pro1} & Canonical representation & Structural form of $f$ \\
		\hline
		Proposition~\protect\ref{Pro2} & One-variable reduction & Growth of $e^h$ \\
		\hline
		Proposition~\protect\ref{Pro3} & Propositions~\protect\ref{Pro1}, \protect\ref{Pro2} & $\rho(f)=\mu(f)$ \\
		\hline
		Proposition~\protect\ref{Pro4} & Proposition~\protect\ref{Pro3} & Deficiency equals one \\
		\hline\hline
	\end{tabularx}
	\vspace{.5cc}
\caption{Summary of results, tools and purposes}
\end{table}



\section{{\bf Shared Value Problems Associated with Shifts}}

Throughout this paper, we work in the complex Euclidean space $\mathbb{C}^n$.  
For $i=1,2,\ldots,n$, we denote
\begin{align*}
	\partial_{z_i}=\frac{\partial}{\partial z_i}, \qquad 
	\partial_{z_i}^{l_i}=\frac{\partial^{l_i}}{\partial z_i^{l_i}},
\end{align*}
and for a multi-index $I=(i_1,i_2,\ldots,i_n)\in \mathbb{Z}_+^n$, we define
\begin{align*}
	\partial^{I}
	=\frac{\partial^{|I|}}{\partial z_1^{i_1}\partial z_2^{i_2}\cdots \partial z_n^{i_n}},
\end{align*}
where $|I|=\sum_{j=1}^{n} i_j$ and $l_i\in \mathbb{Z}^+$.
For a subset $E\subset [0,\infty)$, we introduce the linear measure
\[
m(E):=\int_E dt,
\]
the logarithmic measure
\[
l(E):=\int_{E\cap [1,\infty)} \frac{dt}{t},
\]
and the upper density
\begin{align*}
	\overline{\mathrm{dens}}\;E
	=\limsup_{r\to\infty}\frac{1}{r}\int_{E\cap [1,r]} dt.
\end{align*}
It follows that if $l(E)<+\infty$ or $m(E)<+\infty$, then $E$ has zero upper density.

Let
\[
\mathcal{S}(f)
=\bigl\{ g:\mathbb{C}^n\to\mathbb{P}^1 \text{ meromorphic } :
T(r,g)=o\!\left(T(r,f)\right) \bigr\},
\]
where $r\to\infty$ outside a possible exceptional set of finite linear measure.

\smallskip
Let $f$, $g$, and $a$ be meromorphic functions on $\mathbb{C}^n$. Then there exist pairs of entire functions $(f_1,f_2)$, $(g_1,g_2)$, and $(a_1,a_2)$, each coprime at every point in $\mathbb{C}^n$, such that
\[
f=\frac{f_2}{f_1}, \qquad
g=\frac{g_2}{g_1}, \qquad
a=\frac{a_2}{a_1}.
\]
We say that $f$ and $g$ \emph{share the value $a$ counting multiplicities (CM)} if
\[
\mu^0_{a_1f_2-a_2f_1}
=\mu^0_{a_1g_2-a_2g_1}
\quad (a\not\equiv\infty),\;\;\text{and}\;\;
\mu^0_{f_1}=\mu^0_{g_1}
\quad (a=\infty).
\]
Similarly, $f$ and $g$ are said to \emph{share the value $a$ ignoring multiplicities (IM)} if
\[
\mu^0_{a_1f_2-a_2f_1,1}
=\mu^0_{a_1g_2-a_2g_1,1}
\quad (a\not\equiv\infty),\;\;\text{and}\;\;
\mu^0_{f_1,1}=\mu^0_{g_1,1}
\quad (a=\infty).
\]

\vspace{1.2mm}

Let $f$ be a meromorphic function in $\mathbb{C}^n$. For $z=(z_1,z_2,\ldots,z_n)$ and
$c=(c_1,c_2,\ldots,c_n)\in\mathbb{C}^n$, we write $z+c=(z_1+c_1,z_2+c_2,\ldots,z_n+c_n).$
The shift of $f$ by $c$ is defined as $f(z+c)=f(z_1+c_1,z_2+c_2,\ldots,z_n+c_n)$,
and the corresponding difference operator is given by
\[
\Delta_c f(z)=f(z+c)-f(z).
\]
Inductively, for $m\ge2$, we define the higher-order difference operator by
\begin{align*}
	\Delta_c^m f(z)=\Delta_c^{m-1}f(z+c)-\Delta_c^{m-1}f(z)=\Delta_c\circ\Delta_c^{m-1}f(z),
\end{align*}
where $m\in\mathbb{N}$.

\smallskip
%The classical Nevanlinna value distribution theory, introduced in the 1920s by Rolf Nevanlinna, studies the manner in which meromorphic functions of a single complex variable assume their values. It provides profound insights into the growth, uniqueness, and value-sharing properties of meromorphic functions and has become a cornerstone of complex analysis.

With the development of complex analysis in higher dimensions, mathematicians sought to extend Nevanlinna's foundational results to functions of several complex variables. This endeavor led to the emergence of Nevanlinna theory in several complex variables, which investigates the growth and value distribution of holomorphic and meromorphic mappings from $\mathbb{C}^n$ into complex manifolds and complex projective spaces. 
In recent years, Nevanlinna theory in several complex variables has evolved into a dynamic and rapidly advancing field of research (see \cite{BM}, \cite{ck}-\cite{CLL}, \cite{PVD2}, \cite{HZ1}, \cite{Korhonen-2012}-\cite{LS1}, \cite{FL1}-\cite{FL}, \cite{Liu-Zhang-2018}, \cite{MD1}, \cite{MDP}, \cite{MS}, \cite{MS1}, \cite{MSP}, \cite{Wang-Liu-2023}). 

A particularly active area concerns the application of value distribution theory to problems involving normal families, partial differential equations (PDEs), partial difference equations, and partial differential-difference equations. 


In 2012, Korhonen \cite{Korhonen-2012} established the first difference analogue of the logarithmic derivative lemma, known as the \emph{logarithmic difference lemma}, for meromorphic functions in $\mathbb{C}^m$ with hyper-order strictly less than $\frac{2}{3}$. Subsequently, Cao and Korhonen \cite{ck} extended this result to the case where the hyper-order is strictly less than $1$. These foundational results play a crucial role in the study of meromorphic solutions of complex difference equations and in the uniqueness theory of difference operators in higher dimensions.

We begin by recalling a fundamental result of Heittokangas et al. \cite{Heittokangas-Korhonen-Laine-Rieppo-Zhang-2009}, which establishes a strong structural behavior of meromorphic functions of small order. 
\begin{theo3A} \cite[Theorem 1]{Heittokangas-Korhonen-Laine-Rieppo-Zhang-2009} Let $f$ be a non-constant meromorphic function in $\mathbb{C}$ such that $\rho(f)<2$ and let $c$ be a non-zero complex number. If $f(z+c)$ and $f(z)$ share $a\in\mathbb{C}$ and $\infty$ CM, then $f(z+c)-a=d(f(z)-a)$ for all $z\in\mathbb{C}$, where $d$ is some non-zero complex number.
\end{theo3A}
The following example demonstrates that the order condition ``$\rho(f)<2$'' of Theorem 3.A is indispensable.
\begin{exm3A}\cite{Heittokangas-Korhonen-Laine-Rieppo-Zhang-2009} Let $f(z)=e^{z^{2}} + 1$, and let $c$ be a non-zero finite complex value. Then it immediately follows that $f(z+c)-1=(f(z)-1)e^{2zc+c^{2}}$ for all $z \in \mathbb{C}$. Moreover, we find that
$N(r,1;f)=0$, and hence $1$ is a Borel exceptional value of $f$.
\end{exm3A}

Motivated by the role of Borel exceptional values, for entire functions of finite order, Huang and Zhang \cite{Huang-Zhang-2018} obtained the following refinement. 
\begin{theo3B}\cite[Theorem 2.3]{Huang-Zhang-2018} Let $f(z)$ be a finite-order transcendental entire function in $\mathbb{C}$ with a Borel exceptional value $\alpha \in \mathbb{C}$, let $c \in \mathbb{C}$, and let $a(z) \,(\not\equiv \alpha) \in \mathcal{S}(f)$.
If $f(z+c)$ and $f(z)$ share $a(z)$ CM, then $f(z+c) \equiv f(z)$.
\end{theo3B}

In the same spirit, Huang and Zhang also established a corresponding result involving higher-order difference operators.
\begin{theo3C}\cite[Theorem 2.6.]{Huang-Zhang-2018} Let $f$ be a transcendental entire function in $\mathbb{C}$ such that $\rho(f)<2$.
If $\Delta_c^{k} f(z)$ and $f(z)$ share $0$ CM, where $k \in \mathbb{N}$ and
$c\in\mathbb{C}\backslash \{0\}$ are such that $\Delta_c^{k} f(z) \not\equiv 0$, then 
$\Delta_c^{k} f(z) \equiv df(z)$ for all $z\in\mathbb{C}$, where $d$ is some non-zero complex number.
\end{theo3C}

Recently, Majumder and Das \cite{MD1} extended Theorem 3.A to functions of several complex variables.
\begin{theo3D}\cite[Thorem 1.4]{MD1} Let $f$ be a non-constant meromorphic function in $\mathbb{C}^n$ such that $\rho(f)<2$ and let $c\in\mathbb{C}^n\backslash \{0\}$. If $f(z+c)$ and $f(z)$ share $a\in\mathbb{C}$ and $\infty$ CM, then $f(z+c)-a=d(f(z)-a)$ for all $z\in\mathbb{C}^n$, where $d$ is some non-zero complex number.
\end{theo3D}

The following example confirms that the condition ``$\rho(f)<2$''remains necessary in higher dimensions.
\begin{exm}\cite{MD1}
Let $f(z)=f(z_1,z_2)=e^{z_1^2+z_2^2}+1$ and $c=(c_1,c_2)=(1,0)$. Clearly $f(z+c)$ and $f(z)$ share $1$ and $\infty$ CM, but $f(z+c)-1=e^{2z_1+1}(f(z)-1)$.
\end{exm}

We now present our first main contribution of this paper.
\begin{theo}\label{t2.2} Let $f$ be a transcendental meromorphic function in $\mathbb{C}^n$ such that $\rho(f)<2$.
If $\Delta_c^{k} f(z)$ and $f(z)$ share $0$ and $\infty$ CM, where $k \in \mathbb{N}$ and
$c\in\mathbb{C}^n \backslash \{0\}$ are such that $\Delta_c^{k} f(z) \not\equiv 0$, then $\rho(f)\geq 1$ and
$\Delta_c^{k} f(z) \equiv df(z)$ for all $z\in\mathbb{C}^n$, where $d$ is some non-zero complex number.
\end{theo}

\begin{rem} Obviously Theorem \ref{t2.2} extends Theorem 3.C to the case of higher dimensions.
\end{rem}

Following example ensures the existence of conclusion $\Delta_c^{k} f(z) \equiv df(z)$ for all $z\in \mathbb{C}^n$
of Theorem \ref{t2.2}.

\begin{exm} Let $f:\mathbb{C}^n\rightarrow \mathbb{P}^1$ be defined by $f(z)=\frac{e^{\alpha (z_1+z_2+\cdots+z_n)}}{e^{\beta (z_1+z_2+\cdots+z_n)}+1}$, where $\alpha,\;\beta\in \mathbb{C}\backslash \{0\}$ such that $e^{n\alpha}=3$, $e^{n\beta}=1$ and $c=(1,1,\ldots,1)$. Note that $\Delta_cf(z)=\frac{2e^{\alpha (z_1+z_2+\cdots+z_n)}}{e^{\beta (z_1+z_2+\cdots+z_n)}+1}$.
Clearly $\Delta_cf(z)=2f(z)$ and so $\Delta_cf$ and $f$ share $0$ and $\infty$ CM.
\end{exm} 
\begin{exm}
	Let $f:\mathbb{C}^n \to \mathbb{P}^1$ be defined by
$f(z)=\frac{(e^{2\pi i \lambda/\mu} - 1) e^{\lambda (z_1+\cdots+z_n)}}{e^{\mu (z_1+\cdots+z_n)} - 1}, \quad \lambda,\mu \in \mathbb{C}\setminus\{0\}, \frac{\lambda}{\mu}\not\in\mathbb{Z}$
	and let $c = \left(\frac{2\pi i}{\mu},0,\dots,0\right) \in \mathbb{C}^n$. Then $\Delta_c^k f(z) = \frac{(e^{2\pi i \lambda/\mu} - 1)^{k+1}\, e^{\lambda (z_1+\cdots+z_n)}}{e^{\mu (z_1+\cdots+z_n)} - 1}.$	Here, $\Delta_c^k f(z) \equiv d\, f(z)$ with $d = (e^{2\pi i \lambda/\mu} -1)^{k+1} \neq 0$.
	Clearly: $\Delta_c^k f$  and $f$ share $0$ and $\infty$ CM and $\rho(f)=1<2$.
\end{exm}

The following examples demonstrate that even when $f(z)$ and $\Delta_c^{k} f(z)$ share a non-zero constant $a$ CM, it does not necessarily follow that
$\Delta_c^{k} f(z) \equiv df(z)$ for all $z\in \mathbb{C}^n$, where $d$ is a non-zero complex number.

\begin{exm} Let $f(z)=e^{z_1+z_2+\cdots+z_n}+e^{-(z_1+z_2+\cdots+z_n)}+2$, $c=\left(\frac{\pi i}{n},\ldots,\frac{\pi i}{n}\right)$. Note that $\Delta_cf(z)=-2(e^{z_1+z_2+\cdots+z_n}+e^{-(z_1+z_2+\cdots+z_n)})$.
Clearly  $\Delta_cf$ and $f$ share $\frac{4}{3}$ and $\infty$ CM, but $\Delta_cf\not\equiv df$.
\end{exm}
 \begin{exm}
 	Let $f(z) = \frac{e^{(z_1+\cdots+z_n)} }{e^{2 (z_1+\cdots+z_n)} +1}+1$, $c=\left(\frac{\pi i}{n},\ldots,\frac{\pi i}{n}\right)$. Note that $\Delta_c^k f(z) = (-1)^{k}2^{k}\\ \frac{e^{(z_1+\cdots+z_n)} }{e^{2 (z_1+\cdots+z_n)} +1}.$	Clearly, $\Delta_c^k f(z)$ and  $f(z)$ share  $a = \frac{1}{1+(-1)^{k+1}2^{-k}}$ and $\infty$ CM but $\Delta_c^kf\not\equiv df$.
 \end{exm}
\begin{theo}\label{t2.3} Let $f$ be a finite-order transcendental entire function in $\mathbb{C}^n$ with a Borel exceptional value $\alpha \in \mathbb{C}$ such that $f(0)\neq \alpha$, $c \in \mathbb{C}^n\backslash \{0\}$ and $a(\not\equiv \alpha)$ be an entire function such that $a\in \mathcal{S}(f)$.
If $f(z+c)$ and $f(z)$ share $a(z)$ CM, then $f(z+c) \equiv f(z)$.
\end{theo}

Following example shows that the condition ``$a(\not\equiv \alpha)$'' in Theorem \ref{t2.3} is necessary.

\begin{exm} Let $f(z)=e^{z_1+z_2+z_3}(z_2+z_3+1)(z_1-z_2+1)+1$ and $c=\left(\frac{\pi \iota}{2},\frac{\pi \iota}{2},-\frac{\pi\iota}{2}\right)$. Note that $f(z+c)=\iota e^{z_1+z_2+z_3}(z_2+z_3+1)(z_1-z_2+1)+1$. Clearly $1$ is a Borel exceptional value of $f$ and $f(0)\neq 1$. It is easy to very that $f(z+c)$ and $f(z)$ share $1$ CM, but $f(z+c)\not\equiv f(z)$.
\end{exm}

The following example demonstrates that Theorem \ref{t2.3} fails to hold when $f$ is a transcendental meromorphic function.
\begin{exm} Let $f:\mathbb{C}^n\rightarrow \mathbb{P}^1$ be defined by 
$f(z)= \frac{e^{(z_1+z_2+\cdots+z_n)}-1}{e^{(z_1+z_2+\cdots+z_n)}+1}$ and let $c=\left(c_1,c_2,\cdots,c_n\right)$ such that $e^{(c_1+c_2+\cdots+c_n)}=-1$. Clearly $1$ is a Borel exceptional value of $f$ and $f(0)\neq 1$. It is easy to very that $f(z+c)$ and $f(z)$ share $-1$ CM, but $f(z+c)\not\equiv f(z)$.
\end{exm}

Following example shows that the condition ``$\rho(f)<+\infty$'' in Theorem \ref{t2.3} is necessary.
\begin{exm} Let $f:\mathbb{C}^n\rightarrow \mathbb{P}^1$ be defined by $\displaystyle f(z)=a\left(1-e^{-\sin(z_1+z_2+\cdots+z_n)}\right)$
and $c=(c_1,c_2,\ldots,c_n)=(r_1\pi,r_2\pi,\ldots,r_n\pi)$, where $\sum_{i=1}^nr_i$ is an odd positive integer, a is a non-zero constant. Note that
$\displaystyle f(z+c)=a\left(1-e^{\sin(z_1+z_2+\cdots+z_n)}\right).$
 Clearly $a$ is a Borel exceptional value of $f$ and $f(0)\neq a $. It is easy to very that $f(z+c)$ and $f(z)$ share $2a$ CM, but $f(z+c)\not\equiv f(z)$.
\end{exm}

Next we consider complex partial differential-difference equation
\begin{align}\label{Eq3.1}
\frac{\partial f(z)}{\partial z_j}=f(z+c),
\end{align}
where $k\in\mathbb{Z}[1,n]$ and $c\in\mathbb{C}^n\backslash \{0\}$.
The solutions of (\ref{Eq3.1}) exist, for example, $f(z)=e^{z_1+z_2+\ldots+z_n}$ is a solution of (\ref{Eq3.1}), where
$c =(2k\pi \iota,2k\pi \iota,\ldots,2k\pi \iota)$, $k\in\mathbb{Z}$ and $f(z)=\sin (z_1+z_2+\ldots+z_n)$ or $f(z)=\cos (z_1+z_2+\ldots+z_n)$ are also solutions of (\ref{Eq3.1}) for suitable $c\in\mathbb{C}^n$. Obviously, the equation
(\ref{Eq3.1}) has no rational solutions in $\mathbb{C}^n$.

\begin{theo}\label{t2.4} The order of meromorphic solutions in $\mathbb{C}^n$ of (\ref{Eq3.1}) must satisfy $\rho(f)\geq 1$.
\end{theo}

\begin{theo}\label{t2.5} Let $f$ be an entire solution in $\mathbb{C}^n$ of the equation (\ref{Eq3.1}) such that $f(0)\neq 0$ and $\rho(f)$ be finite. If $0$ is a Borel exceptional value of $f$, then $\rho(f)=1$. Furthermore, one of the following cases holds:
\begin{enumerate}
\item[(i)] $f(z)=(b_1z_1+\ldots+b_nz_n+b_0)e^{a_1z_1+\ldots+a_nz_n+a_0}$, where $a_0,a_1,\ldots,a_n,b_0,b_1,\ldots,b_n\in\mathbb{C}$ such that $(a_1,\ldots,a_n)\neq (0,\ldots,0)$, $b_0\neq 0$ and $a_j=e^{a_1c_1+a_2c_2+\ldots+a_nc_n}$;
\item[(ii)] $f(z)=g(z)e^{a_1z_1+\ldots+a_nz_n+a_0}$, where $a_0,a_1,\ldots,a_n\in\mathbb{C}$ such that $(a_1,\ldots,a_n)\neq (0,\ldots,0)$, $a_j=e^{a_1c_1+a_2c_2+\ldots+a_nc_n}$ and $g(z)$ is a transcendental entire function in $\mathbb{C}^n$ such that $\rho(g)<1$ and $\frac{\partial g(z)}{\partial z_j}=e^{a_1c_1+a_2c_2+\ldots+a_nc_n}(g(z+c)-g(z))$.
\end{enumerate}
\end{theo}

\begin{rem} Obviously Theorem \ref{t2.5} extends Theorem 3.5 \cite{Liu-Dong-2014} to the case of higher dimensions.
\end{rem}

\subsection{{\bf Auxiliary lemmas}}
\begin{lem}\label{L2.1}\cite[Theorem 2.4]{Cao-Xu-2020} Let $f$ be a non-constant meromorphic function on $\mathbb{C}^n$ and let $c \in \mathbb{C}^n\setminus \{0\} $. If $f$ is of finite order, then 
\begin{align*}
 m\left(r,\frac{f(z+c)}{f(z)}\right)+m\left(r, \frac{f(z)}{f(z+c)}\right)=O\left(r^{\rho(f)-1+\varepsilon}\right)
 \end{align*}
holds for any $\varepsilon>0$.
\end{lem}

\begin{lem}\label{L2.2}\cite[Corollary 3.2]{Cao-Xu-2020} Let $A_0$, $\ldots$, $A_m$ be entire functions in $\mathbb{C}^n$ such that there exists an integer $k \in \{0, \ldots, m\}$ satisfying 
\begin{align*}
\rho(A_k)>\max \{\rho(A_j): 0 \le j \le m, j \ne k\}.
\end{align*}
If $f$ is a nontrivial entire solution of linear partial difference equation 
\begin{align*}
A_m(z)f(z+c_m)+\ldots+A_1(z)f(z+c_1)+A_0f(z)=0,
\end{align*}
 where $c_1,\ldots,c_m$ are distinct values of $\mathbb{C}^n\setminus\{0\}$, then we have $\rho(f)\ge \rho(A_k)+1$.
\end{lem}

\begin{lem}\label{L2.3}\cite[Theorem 2.2]{Cao-Xu-2020} Let $f$ be a non-constant meromorphic function on $\mathbb{C}^n$ with 
\begin{align*}
\lim\limits_{r\rightarrow \infty}\sup \frac{\log T(r,f)}{r}=0,
\end{align*}
then $\parallel\;T(r,f(z+c))=T(r,f)+o(T(r,f))$ holds for any constant $c\in\mathbb{C}^n$.
\end{lem}

\begin{lem}\label{L2.4} \cite[Theorem 1.26]{Hu-Li-Yang-2003} Let $f:\mathbb{C}^n\to\mathbb{P}^1$ be a non-constant meromorphic function. Assume that 
$R(z, w)=\frac{A(z, w)}{B(z, w)}$. Then
\beas \parallel\;T\left(r, R_f\right)=\max \{p, q\} T(r, f)+O\Big(\sideset{}{_{j=0}^{p}}{\sum}T(r, a_j)+\sideset{}{_{j=0}^{q}}{\sum}T(r, b_j)\Big),\eeas
where $R_f(z)=R(z, f(z))$ and two coprime polynomials $A(z, w)$ and $B(z,w)$ are given
respectively $A(z,w)=\sum_{j=0}^p a_j(z)w^j$ and $B(z,w)=\sum_{j=0}^q b_j(z)w^j$.
\end{lem}

\begin{lem}\label{L2.5}\cite[Corollary 4.5]{Berenstein-Chang-Li-1994} Suppose $a_0 ,a_1,\ldots,a_m\; (m\geq 1)$ are meromorphic functions on $\mathbb{C}^n$ and $g_0,g_1,\ldots,g_m$ are entire functions on $\mathbb{C}^n$ such that $g_j-g_k$ are not constants for $0\leq j<k\leq m$. If the conditions 
\begin{align*}
\sum\limits_{j=0}^m a_je^{g_j}=0
\end{align*}
and $\parallel T(r,a_j)=o(T(r))$, $j=0,1,\ldots,n$
hold, where $T(r)=\max_{0\leq k<j\leq n}(T(r,e^{g_j-g_k}))$,
then $a_j\equiv 0$ for all $j=0,1,\ldots,n$.
\end{lem}

\begin{lem}\label{L2.6}\cite{MS1} Let $f$ be a transcendental meromorphic function in $\mathbb{C}^n$ such that $\rho(f)<1$. Let $h>0$. Then there exists an $\varepsilon$-set $E$ in $\mathbb{C}^n$ such that
\begin{align*}
\frac{\partial_{z_i}(f(z+c))}{f(z+c)}\rightarrow 0
\end{align*}
for all $i\in\mathbb{Z}[1,n]$ and 
\begin{align*}
\frac{f(z+c)}{f(z)}\rightarrow 1
\end{align*}
as $||z||\rightarrow \infty$, where $z\in\mathbb{C}^n\backslash  E$ and $||c||\leq h$. Further, $E$ may be chosen so that for large $||z||$, where $z\not\in E$, the function $f(z)$ has no zeros or poles in $||\zeta-z||\leq h$.
\end{lem}

\begin{lem}\label{L2.7}\cite{MS1} Let $f$ be a transcendental meromorphic function in $\mathbb{C}^n$ of finite order $\rho$. 
Then there exists an $\varepsilon$-set $E$ in $\mathbb{C}^n$ such that
\begin{align*}
 \displaystyle \left|\frac{\partial_{z_i}(f(z))}{f(z)}\right|\leq ||z||^{\rho-1+\delta},
 \end{align*}
holds for all large values of $||z||$, where $z\in \mathbb{C}^n\backslash E$ and $\delta>0$ is a given constant. 
\end{lem}


\begin{lem}\label{L2.8} Let $P$ and $Q$ be two polynomials in $\mathbb{C}^n$ such that $P$ is non-constant. If $f$ is a finite order solution in $\mathbb{C}^n$ of the following equation
\begin{align*}
f(z+c)=e^{P(z)}\left(\frac{\partial f(z)}{\partial z_j}+Q(z)f(z)\right)
\end{align*}
where $c\in\mathbb{C}^n\backslash \{0\}$ and $j\in\mathbb{Z}[1,n]$, then $\rho(f)\geq \deg(P)+1$.
\end{lem}
\begin{proof} We have
\begin{align}\label{Eq3.2}
\frac{f(z+c)}{f(z)}=e^{P(z)}\left(\frac{\partial f(z)}{\partial z_j}/f(z)+Q(z)\right).
\end{align}
Now using Lemma 1.37 \cite{Hu-Li-Yang-2003}, we have
\begin{align*}
m\left(\frac{\partial f(z)}{\partial z_j}/f(z)\right)=O(\log r)
\end{align*}
and by Lemma \ref{L2.1}, we obtain
\begin{align*}
m\left(r,f(z+c)/f(z)\right)=O\left(r^{\rho(f)-1+\varepsilon}\right)
 \end{align*}
holds for any $\varepsilon>0$. Since $m(r,Q)=O(\log r)$, from (\ref{Eq3.2}), we deduce that
\begin{align*}
m(r,e^P)=O(r^{\deg(P)})\leq O\left(r^{\rho(f)-1+\varepsilon}\right)+O(\log r),
\end{align*}
which shows that $\rho(f)\geq \deg(P)+1$.
\end{proof}

\subsection{{Proof of Theorem \ref{t2.2}}} By the given condition, we have
\begin{align}\label{t2.2:1.1}
\frac{\Delta_c^k f(z)}{f(z)}=e^{P(z)},
\end{align}
where $P$ is an entire function in $\mathbb{C}^n$. Since $\rho(f)<2$, using Lemma 2.B to (\ref{t2.2:1.1}), we deduce that $P$ is a polynomial in $\mathbb{C}^n$ of degree at most one. If possible, suppose $\deg(P)=1$. Note that
\begin{align}\label{t2.2:1.2}
\Delta_c^{k} f(z)=\sum_{j=0}^{k}(-1)^{k-j}\binom{k}{j}f(z + jc).
\end{align}
Substituting (\ref{t2.2:1.2}) into (\ref{t2.2:1.1}), we get
\begin{align}\label{t2.2:1.3}
f(z + kc)+\sum_{j=1}^{k-1}(-1)^{k-j}\binom{k}{j}f(z + jc)+\bigl[(-1)^k - e^{P(z)}\bigr]f(z)= 0.
\end{align}

Now using Lemma \ref{L2.2} to (\ref{t2.2:1.3}), we deduce that $\rho(f)\ge \rho(e^P)+ 1=2$, which contradicts the assumption that $\rho(f)<2$. Hence $P$ is a constant. If we take $d=e^P$, then from (\ref{t2.2:1.1}), we have 
\begin{align}\label{t2.2:1.4}
\Delta_c^{k} f(z) \equiv df(z)
\end{align}
for all $z\in\mathbb{C}^n$. If possible, suppose $\rho(f)<1$. Now substituting (\ref{t2.2:1.2}) into (\ref{t2.2:1.4}) and then using Lemma \ref{L2.6}, we get
\begin{align*}
d=\frac{\Delta_c^{k} f(z)}{f(z)}=\sum_{j=0}^{k}(-1)^{k-j}\binom{k}{j}\frac{f(z + jc)}{f(z)}\to \sum_{j=0}^{k}(-1)^{k-j}\binom{k}{j}=(1-1)^k=0
\end{align*}
as $||z||\to\infty$ possibly outside an $\varepsilon$-set $E$ in $\mathbb{C}^n$, which is impossible. Hence $\rho(f)\geq 1$.

\subsection{{Proof of Theorem \ref{t2.3}}} By the given condition, we have
\begin{align}\label{t2.3:1.1}
\frac{f(z+c)-a(z)}{f(z)-a(z)}=e^{P(z)},
\end{align}
where $P$ is an entire function in $\mathbb{C}^n$. Using Lemmas \ref{L2.3} and \ref{L2.4} to (\ref{t2.3:1.1}), we deduce that $\rho(e^P)\leq \rho(f)$. Since $\rho(f)$ is finite, using Lemma 2.B, we obtain $\deg(P)=\rho(e^P)\leq \rho(f)$. Note that $\infty$ is a Borel exceptional value of $f$. Since $\alpha\in\mathbb{C}$ is a Borel exceptional value of $f$, by Proposition \ref{Pro1}, we have 
\begin{align}\label{t2.3:1.2}
f(z)=H(z)e^{Q(z)}+\alpha,
\end{align}
where $H(z)$ is a non-zero entire function in $\mathbb{C}^n$ and $Q(z)$ is a polynomial in $\mathbb{C}^n$ such that $\rho(H)<\rho(f)=\rho(e^Q)=\deg(Q)$. Since $f$ has two distinct Borel exceptional values, by Proposition \ref{Pro3}, we have $\rho(f)=\mu(f)$. Then by Proposition 1.A, we have $T(r,H)=o(T(r,f))$. We now conclude from (\ref{t2.3:1.1}) and (\ref{t2.3:1.2}) that
\begin{align}\label{T2.3:1.3}
H(z+c)e^{Q(z+c)}-H(z)e^{P(z)+Q(z)}+(a(z)-\alpha)e^{P(z)}-(a(z)-\alpha)= 0.
\end{align}

First, we suppose that $P(z)$ is a non-constant polynomial. Clearly
\begin{align*}
1 \le \deg(P) \le \rho(f) = \deg(Q).
\end{align*}
We now consider the following two cases.

\medskip
\noindent
\textbf{Case 1.} Let $1 \le \deg(P) < \rho(f)=\deg(Q)$. Equation (\ref{T2.3:1.3}) can be rewritten as
\begin{align}\label{t2.3:1.3}
H(z+c)e^{Q(z+c)-Q(z)}-H(z)e^{P(z)}=(a(z)-\alpha)\bigl(1 - e^{P(z)}\bigr)e^{-Q(z)}.
\end{align}

Using Proposition 1.A, we see that the order of the left-hand side of (\ref{t2.3:1.3}) is less than $\rho(f)$, while the order of the right-hand side of (\ref{t2.3:1.3}) is equal to $\rho(f)$. This yields a contradiction.

\medskip
\noindent
\textbf{Case 2.} Let $1 \le \deg(P)=\rho(f)=\deg(Q)=m$. Suppose
\begin{align*}
P(z)=\sum\limits_{j=0}^m P_j(z)\quad \text{and}\quad Q(z)=\sum\limits_{j=0}^m Q_j(z),
\end{align*}
where $P_j(z)$ and $Q_j(z)$ are homogeneous polynomials in $\mathbb{C}^n$ of degree $j\;(0\le j\le m)$ such that $P_m(z)$ and $Q_m(z)$ are non-constant.
Now we consider the following three sub-cases.

\medskip
\noindent
\textbf{Sub-case 2.1.} Let $P_m(z)-Q_m(z)\equiv 0$. Equation (\ref{t2.3:1.3}) can be rewritten as
\begin{align}\label{t2.3:1.4}
H_{11}(z)e^{-Q(z)} + H_{12}(z)e^{P(z)} + H_{13}(z)e^{Q_{0}(z)}=0,
\end{align}
where $Q_{0}(z)\equiv 0$ and
\begin{align*}
H_{11}(z) = \alpha - a(z),\; H_{12}(z) = -H(z),\;
H_{13}(z)= H(z+c)e^{Q(z+c)-Q(z)} + (a(z)-\alpha)e^{P(z)-Q(z)}.
\end{align*}
Note that
\begin{align*}
\deg(-Q-P)= \deg(-Q-Q_{0})= \deg(P-Q_{0})=m,
\end{align*}
and $T(r,a)=o(T(r,e^{Q})).$
Thus, for all $j = 1,2,3$, we have
\begin{align*}
T(r,H_{1j})= o\left(T\left(r,e^{-Q-P}\right)\right),\;\; T(r,H_{1j})= o\left(T\left(r,e^{-Q-Q_{0}}\right)\right)
\end{align*}
and $T(r,H_{1j})= o\left(T\left(r,e^{P-Q_{0}}\right)\right).$
Now using Lemma \ref{L2.5} to (\ref{t2.3:1.4}), we deduce that $H_{1j}\equiv 0$ for $j = 1,2,3$, which is a contradiction.

\medskip
\noindent
\textbf{Sub-case 2.2.} Let $P_m(z)+Q_m(z)\equiv 0$. Equation (\ref{t2.3:1.3}) can be rewritten as
\begin{align*}
H_{21}(z)e^{Q(z+c)} + H_{22}(z)e^{P(z)} + H_{23}(z)e^{Q_{0}(z)} = 0,
\end{align*}
where $Q_{0}\equiv 0$ and
\begin{align*}
H_{21}(z) = H(z+c),\;\;H_{22}(z) = a(z) - \alpha\;\text{and}\; H_{23}(z) = -H(z)e^{Q(z)+P(z)} - (a(z)-\alpha).
\end{align*}

Now proceeding in the same way as done in the proof Sub-case~2.1, we deduce that $H_{2j}\equiv 0$ for $j = 1,2,3$, which is a contradiction.

\medskip
\noindent
\textbf{Sub-case 2.3.} Let $P_m(z)\pm Q_m(z)\not\equiv 0$. Equation (\ref{t2.3:1.3}) can be rewritten as
\begin{align*}
H_{31}(z)e^{Q(z+c)} + H_{32}(z)e^{Q(z)+P(z)}+ H_{33}(z)e^{P(z)} + H_{34}(z)e^{Q_{0}(z)}=0,
\end{align*}
where $Q_{0}\equiv 0$ and
\begin{align*}
H_{31}(z) = H(z+c),\;\;H_{32}(z)= -H(z),\;\;H_{33}(z)= a(z) - \alpha\;\text{and}\;H_{34}(z) = \alpha - a(z).
\end{align*}
Now proceeding in the same way as done in the proof Sub-case~2.1, we deduce that $H_{3j}\equiv 0$ for $j=1,2,3,4$, which is a contradiction.
\vspace{1.2mm}

Next we suppose that $P$ is a constant. Let $d=e^{P}$. If possible, suppose $d\neq 1$. Then equation (\ref{t2.3:1.3}) becomes
\begin{align}\label{t2.3:1.5}
H(z+c)e^{Q(z+c)-Q(z)}-dH(z)=(a(z)-\alpha)(1-d)e^{-Q(z)}.
\end{align}
Note that
\begin{align*}
\rho\!\left((a(z)-\alpha)(1-d)e^{-H(z)}\right)=m= \rho(f)
\end{align*}
and
\begin{align*}
\rho\!\left(H(z+c)e^{Q(z+c)-Q(z)}-dH(z)\right)<m= \rho(f).
\end{align*}
Therefore from (\ref{t2.3:1.5}), we get a contradiction. Hence $d=1$ and so from (\ref{t2.3:1.1}), we have reduces to $f(z+c) \equiv f(z)$.


\subsection{{Proof of Theorem \ref{t2.4}}} It is easy to verify that the equation
(\ref{Eq3.1}) has no rational solutions in $\mathbb{C}^n$. If possible, suppose $f$ is a transcendental meromorphic solution of the equation (\ref{Eq3.1}) such that $\rho(f)<1$. Then By Lemmas \ref{L2.6} and \ref{L2.7}, we have respectively
\begin{align}\label{t2.4:1.1}
\displaystyle \frac{f(z+c)}{f(z)}\rightarrow 1 \;\text{and}\; \left|\frac{\partial_{z_i}(f(z))}{f(z)}\right|=o(1),
\end{align}
for large values of $||z||$, where $z\in \mathbb{C}^m\backslash E$ and $i\in\mathbb{Z}[1,m]$, where $E$ is an $\varepsilon$-set in $\mathbb{C}^n$. Now using (\ref{t2.4:1.1}) to (\ref{Eq3.1}), we immediately get a contradiction. Hence $\rho(f)\geq 1$.


\subsection{{Proof of Theorem \ref{t2.5}}} Note that $0$ and $\infty$ are the Borel exceptional values of $f$.
By Proposition \ref{Pro1}, we have $f(z)=g(z)e^{Q(z)}$, where $g(z)$ is a non-zero entire function in $\mathbb{C}^n$ such that $g(0)\neq 0$ and $Q(z)$ is a polynomial in $\mathbb{C}^n$ such that $\rho(g)<\rho(f)=\rho(e^Q)=\deg(Q)$. Also by Proposition \ref{Pro3}, we have $\rho(f)=\mu(f)\in\mathbb{N}$. Then by Proposition 1.A, we have $T(r,g)=o(T(r,f))$. Now from (\ref{Eq3.1}), we have
\begin{align}\label{t2.5:1.1}
\left(\frac{\partial g(z)}{\partial z_j}+g(z)\frac{\partial Q(z)}{\partial z_j}\right)e^{Q(z)-Q(z+c)}=g(z+c).
\end{align}
Note that $\deg(Q(z)-Q(z+c))=\deg(Q(z))-1$. If $Q(z)-Q(z+c)$ is non-constant, then using lemma \ref{L2.7} to (\ref{t2.5:1.1}), we deduce that $\rho(g)\geq \deg(Q(z)-Q(z+c))+1=\deg(Q(z))$, which contradicts the fact that $\rho(g)<\deg(Q)$. Hence $Q(z)-Q(z+c)$ is a constant and so $\deg(Q)=1$. Let $Q(z)=a_1z_1+a_2z_2+\ldots+a_nz_n+b$, where $(a_1,a_2,\ldots,a_n)\neq (0,0,\ldots,0)$ and $b\in\mathbb{C}$. Clearly $Q(z+c)-Q(z)=a_1c_1+a_2c_2+\ldots+a_nc_n$ and $\rho(g)<1$. Now from (\ref{t2.5:1.1}), we get
\begin{align}\label{t2.5:1.2}
\frac{\partial g(z)}{\partial z_j}+a_jg(z)=e^{a_1c_1+a_2c_2+\ldots+a_nc_n}g(z+c).
\end{align}
Using (\ref{t2.4:1.1}) to (\ref{t2.5:1.2}), we deduce that $a_j=e^{a_1c_1+a_2c_2+\ldots+a_nc_n}$ and so from (\ref{t2.5:1.2}), we get
\begin{align}\label{t2.5:1.3}
\frac{\partial g(z)}{\partial z_j}=a_j(g(z+c)-g(z)).
\end{align}
If $g$ is a polynomial in $\mathbb{C}^n$, then from (\ref{t2.5:1.3}) it is easy to deduce that $\deg(g)\leq 1$. Let $g(z)=b_1z_1+b_2z_2+\ldots+b_nz_n+b_0$, where $b_0,b_1,b_2,\ldots,b_n\in\mathbb{C}$. Since $g(0)\neq 0$, it follows that $b_0\neq 0$.



\vspace{0.1in}
{\bf Compliance of Ethical Standards:}\par

{\bf Conflict of Interest.} The authors declare that there is no conflict of interest regarding the publication of this paper.\par

{\bf Data availability statement.} Data sharing not applicable to this article as no data sets were generated or analysed during the current study.


\begin{thebibliography}{99}
\bibitem{BM} {\sc A. Banerjee} and {\sc S. Majumder}, Analytic perspectives on characterizing unique range set of meromorphic functions in several complex variables, \textit{Bull. Korean Math. Soc.}, (accepted for publication).

\bibitem{Berenstein-Chang-Li-1994} {\sc C. Berenstein}, {\sc D. C. Chang} and {\sc B. Q. Li}, A note on wronskians and linear dependence of entire functions in $\mathbb{C}^n$, \textit{Complex Var. Theory Appl.,} \textbf{24} (1-2) (1994), 131-144.

\bibitem{ck} {\sc T. B. Cao} and {\sc R. J. Korhonen}, A new version of the second main theorem for meromorphic mappings intersecting hyperplanes in several complex variables, \textit{J. Math. Anal. Appl.}, \textbf{444} (2) (2016), 1114-1132.

\bibitem{CK1} {\sc T. B. Cao} and {\sc R. J. Korhonen}, Growth of meromorphic solutions and delay partial differential equations, \textit{Israel J. Math.}, (2025), 1-27, DOI:10.1007/s11856-025-2756-8.

\bibitem{Cao-Xu-2020} {\sc T. B. Cao} and {\sc L. Xu}, Logarithmic difference lemma in several complex variables and partial difference equations, \textit{Ann. Mat. Pura Appl.}, \textbf{199} (2020), 767-794.

\bibitem{CLL} {\sc Q. Cheng}, {\sc Y. Li} and {\sc Z. Liu}, On meromorphic solutions of nonlinear partial differential-difference equations of first order in several complex variables, \textit{Bull. Korean Math. Soc.}, \textbf{60} (2) (2023), 425-441.

\bibitem{CY1} {\sc W. Cherry} and  {\sc Z. Ye}, Nevanlinna's theory of value distribution second main theorem and its error terms. Springer, Berlin (2001).

\bibitem{PVD2} {\sc P. V. Dovbush}, Zalcman-Pang's lemma in $C^N$, \textit{Complex Var. Elliptic Equ.}, \textbf{66} (12) (2021), 1991-1997.

\bibitem{HZ1} {\sc W. Hao} and {\sc Q. Zhang}, Meromorphic solutions of a class of nonlinear partial differential equations, \textit{Indian J. Pure Appl. Math.}, 2025: 1-11, doi.org/10.1007/s13226-025-00779-5.

\bibitem{Heittokangas-Korhonen-Laine-Rieppo-Zhang-2009} {\sc J. Heittokangas}, {\sc R. Korhonen}, {\sc I. Laine}, {\sc J. Rieppo}
and {\sc J. L. Zhang}, Value sharing results for shifts of meromorphic functions, and sufficient conditions for periodicity, \textit{J. Math. Anal. Appl.}, \textbf{355} (1) (2009), 352-363.

\bibitem{AH} {\sc A. Hinkkanen}, A sharp form of Nevanlinna's second fundamental theorem, \textit{Invent. Math.}, \textbf{108} (1992), 549-574.

\bibitem{LH1} {\sc L. H\"{o}rmander}, An introduction to complex analysis in several variables, Van Nostrand. Princeton, N. J. 1973.

\bibitem{Hu-Yang-1996} {\sc P. C. Hu} and {\sc C. C. Yang}, Uniqueness of meromorphic functions on $\mathbb{C}^m$, \textit{Complex Var. Elliptic Equ.,} \textbf{30} (1996), 235-270.

\bibitem{Hu-Li-Yang-2003} {\sc P. C. Hu}, {\sc P. Li} and {\sc C. C. Yang}, Unicity of Meromorphic Mappings. Springer, New York (2003).

\bibitem{ps1} {\sc P. C. Hu} and {\sc C. C. Yang}, The Tumura-Clunie theorem in several complex variables, \textit{Bull. Aust. Math. Soc.}, \textbf{90} (2014), 444-456.

\bibitem{Huang-Zhang-2018} {\sc Z. B. Huang} and {\sc R. R. Zhang}, Uniqueness of the differences of meromorphic functions, \textit{Anal. Math.,} \textbf{44} (2018), 461-473. 

\bibitem{Korhonen-2012} {\sc R. J. Korhonen}, A difference Picard theorem for meromorphic functions of several variables, \textit{Comput. Methods Funct. Theo.,}  \textbf{12} (2012) (1), 343-361.

\bibitem{LY11} {\sc B. Q. Li} and {\sc L. Yang}, Picard type theorems and entire solutions of certain nonlinear partial differential equations, \textit{J. Geom. Anal.}, (2025) 35:234, https://doi.org/10.1007/s12220-025-02067-4.

\bibitem{LS1} {\sc Y. Li} and {\sc H. Sun}, A note on unicity of meromorphic functions in several variables, \textit{J. Korean Math. Soc.}, \textbf{60} (4) (2023), 859-876.

\bibitem{Liu-Dong-2014} {\sc K. Liu} and {\sc X. Dong}, Some results related to complex differential-difference equations of certain types, \textit{Bull. Korean Math. Soc.}, \textbf{51} (95) (2014), 1453-1467.

\bibitem{FL1} {\sc F. L\"{u}}, Theorems of Picard type for meromorphic function of several complex variables, \textit{Complex Var. Elliptic Equ.}, \textbf{58} (8) (2013), 1085-1092.

\bibitem{LB} {\sc F. L\"{u}} and {\sc W. Bi}, On entire solutions of certain partial differential equations, \textit{J. Math. Anal. Appl.}, \textbf{516} (1) (2022), 126476.

\bibitem{FL} {\sc F. L\"{u}}, On meromorphic solutions of certain partial differential equations, \textit{Canadian Math. Bull.}, 2025:1-15, doi:10.4153/S0008439525000347.

\bibitem{Liu-Zhang-2018} {\sc Z. Liu} and {\sc Q. Zhang}, Difference uniqueness theorems on meromorphic functions in several variables, \textit{Turk. J. Math.,} \textbf{42} (2018), 2481-2505.

\bibitem{MD1} {\sc S. Majumder} and {\sc P. Das}, Periodic behavior of meromorphic functions sharing values with their shifts in several complex variables, \textit{Indian J. Pure Appl. Math.} https://doi.org/10.1007/s13226-025-00778-6.

\bibitem{MDP} {\sc S. Majumder}, {\sc P. Das} and {\sc D. Pramanik}, Sufficient condition for entire solution of a certain type of partial differential equation in $\mathbb{C}^m$, \textit{J. Contemp. Math. Anal.}, \textbf{60} (5) (2025), 378-395.

\bibitem{MS} {\sc S. Majumder} and {\sc N. Sarkar}, Periodic behavior of meromorphic functions sharing values with their difference operators in several complex variables, \textit{Indian J. Pure Appl. Math.} https://doi.org/10.1007/s13226-025-00916-0.

\bibitem{MS1} {\sc S. Majumder} and {\sc N. Sarkar}, Bergweiler-Langley lemmas in several complex variables, 
\textit{Bull. Belgian Math. Soc.} (accepted for publication).

\bibitem{MSP} {\sc S. Majumder}, {\sc N. Sarkar} and {\sc D. Pramanik}, Solutions of complex Fermat-type difference equations in several variables, \textit{Houston J. Math.} (accepted for publication).


\bibitem{Nishino-2001} {\sc T. Nishino}, Function theory in several complex variables, Vol. 193, American Math. Soc., 2001.

\bibitem{Ronkin-1968} {\sc L. I. Ronkin}, An analog of the canonical product for entire functions of several complex variables, \textit{Trans. Moscow Math. Soc.}, \textbf{18} (1968), 117-160.

\bibitem{Stoll-1953} {\sc W. Stoll}, Ganze Funktionen endlicher Ordnung mit gegebenan Nullstellenfl\"{a}chen, \textit{Math. Z.}, \textbf{57} (1953), 211-237.

\bibitem{WS4} {\sc W. Stoll}, Normal families of non-negative divisors, \textit{Math Z.}, \textbf{84} (1964), 154-218.

\bibitem{Stoll-1974} {\sc W. Stoll}, Holomorphic functions of finite order in several complex variables, Conference Board of the Mathematical Sciences, Regional Conference Series in Mathematics 21, Amer. Math. Soc., 1974.

\bibitem{Stoll-1983} {\sc W. Stoll}, Value distribution and the lemma of the logarithmic derivative on polydiscs, \textit{Internat. J. Math. \& Math. Sci.}, \textbf{6} (4) (1983), 617-669.

\bibitem{Stoll-1985} {\sc W. Stoll}, Value distribution theory of meromorphic maps, \textit{Aspects of Math.} E7 (1985), pp. 347, Vieweg-Verlag.

\bibitem{Wang-Liu-2023} {\sc M. Wang} and {\sc Z. Liu}, Meromorphic functions sharing small functions with their difference polynomial in several variables, \textit{Hacet. J. Math. Stat.}, \textbf{52} (1) (2023), 73-90.

\bibitem{Yang-Yi-2003} {\sc C. C. Yang} and {\sc H. X. Yi}, Uniqueness theory of meromorphic functions, Kluwer Academic Publishers, Dordrecht/Boston/London, 2003.

\bibitem{ZY1} {\sc Z. Ye}, On Nevanlinna's second main theorem in projective space, \textit{Invent. Math.}, \textbf{122} (1995), 475-507.

\end{thebibliography}
\end{document}
